\thispagestyle{empty}
\begin{center}
    \LARGE\textbf{\scshape{\OlympString{tournamelowercase}}}%
\end{center}
%
\vspace{2pt}%
%
\begin{center}
    \large{\OlympString{touronedate}}
\end{center}
\vspace{2pt}
\begin{center}
    \large\textbf{Алдымен төмендегі ақпаратты оқып шығыңыз:}
\end{center}
\vspace{4pt}
\begin{itemize}[leftmargin=*]
    \item Теориялық тур \OlympString{problemcount_instructions} тұрады.
    \item Шешіміңізге қажетті математикалық кеңестер мен физикалық тұрақтылар тізімі тиісті есептердің шартында берілген.
    \item Бұл жинақта сіздің тіліңізде, титулдық парақ пен сіз қазір оқып отырған нұсқаулық парағын \textit{қоспағанда}, барлығы \pageref{LastPage} бет бар.
    \item Есептерді шығару үшін сізге кеңсе жабдықтары мен инженерлік калькулятор қажет.
    \item Сізге шимай жазба үшін таза қағаз және \textbf{\textit{Қатысушының шешімдерін толтыруға арналған өрістер}} (таза парақтар) беріледі. Қағаздар жетіспеген жағдайда қолыңызды көтеріңіз, жақындаған жауапты адамға (бақылаушыға) мәселеңізді айтыңыз --- сізге қосымша парақтар беріледі.
    \item Жұмысыңыз \textbf{\textit{таза парақтар}} бойынша бағаланады; шимай қағаздардағы жазбалар бағаланбайды. Жалпы жағдайда сіз теңдеулер, сандық және әріптік белгілеулер, суреттер мен сызбалар (графиктер) қолдануыңыз қажет; есеп шартында сұралмаған жағдайда сөздік сипаттамалар бағаланбайды.
    \item Есепті шығару барысында есептің нөмерін нақты көрсетіңіз. Әр жаңа есепті \textbf{\textit{таза парақтың}} жаңа бетінен бастаған жөн.
    \item \textbf{\textit{Таза парақтардың}} тек қана алдыңғы бетін қолданыңыз және шекараның сыртынан шықпауға тырысыңыз. \textbf{\textit{Таза парақтардың}} жоғары бөлігінде \textit{беттің реттік нөмерін} жазуыңыз керек. Егер де сіз белгілі бір таза парақтардың бағаланғанын қаламасаңыз, оларды толықтай үлкен айқышпен (крестпен) сызып тастаңыз және толық қолданылған парақтар санына қоспаңыз.
    \item \textbf{Таза парақтарда өзіңіз туралы ақпаратты (аты-жөніңіз, мектеп, т.б.) көрсетпеңіз.} Тур аяқталғаннан кейін жұмыстарыңызды ұйымдастырушылар шифрлайды.
    \item Қосымша кеңестер:
    \begin{itemize}[label=$\circ$]
        \item Есептің шартын толықтай аяғына дейін оқып, оны түсінгеніңізді нақтылаңыз. Кейде есептердің кейбір пункттерін өткізіп, келесі пункттерін дұрыс шығаруға болады.
        \item Есепті жауабына дейін шығара алмасаңыз да бүкіл қажетті теңдеулерді толықтай нақты жазыңыз.
        \item Есептерді қайта тексеру мүмкіндігіне, сондай-ақ жауап пен аралық формулалардың өлшемділігін (размерность) тексеруге немқұрайлы қарамаңыз --- бір түкке тұрмайтын қате үшін талай ұпайдан айырылып қалуға болады!
        \item Есепті шығару барысында шешіміңізді бірден \textbf{\textit{таза парақтарда}} жазуға, ал шимай қағазды тек қана аралық есептеулерді жүргізу үшін, ауқымды теңдеулер жүйесін шешу және алгебралық ықшамдау үшін қолдануға тырысыңыз. Шимай қағаздан қайта көшіріп жазуға артық уақыт жұмсамай, барлық шешімдеріңізді таза параққа түсіруге тырысуыңыз керек.
        \item Басқа ешқандай идея жоқ болып көрінсе де, бөлінген уақытты толықтай қолданыңыз. Жиі жағдайда дұрыс шешім әдісі ұзақ ойланғаннан кейін келеді. Сонымен қатар, сіз тапқан шаманың дәл сол сұралған шама екеніне көз жеткізу үшін шешімді жазғаннан кейін есептің шартын қайта оқып шығу маңызды.
    \end{itemize}
\end{itemize}
\vspace{4pt}
\begin{center}
    \textsc{Сізге сәттілік тілейміз!}
\end{center}
%
\begin{center}
    \textit{\OlympString{duration}}
\end{center}

\newpage